\section{Implementação do sistema}
\label{sec:implementacao-real}
A implementação foi inspecionada e executada em ambiente isolado. Esta seção apresenta os componentes, os testes e as evidências obtidas em 7 de setembro de 2026, no horário de Brasília. Os arquivos de execução utilizam UTC; por isso, parte dos registros está datada de 8 de setembro. As verificações demonstram os comportamentos observados nas condições testadas, sem substituir o experimento comparativo do Capítulo~\ref{cap:metodo}.

\subsection{Tecnologias utilizadas}
A imagem da aplicação utiliza Python 3.12-slim; a execução de integração registrou Python 3.12.14. As dependências fixadas são FastAPI 0.116.1, Uvicorn 0.35.0, Pydantic 2.11.7, SQLAlchemy 2.0.42, psycopg 3.2.9, confluent-kafka 2.11.0, pytest 8.4.1 e httpx 0.28.1. FastAPI fornece as rotas, Pydantic valida o contrato, SQLAlchemy e psycopg executam transações e confluent-kafka fornece produtor e consumidor.

PostgreSQL 16-alpine, Kafka 3.9.1 e k6 1.4.0 foram executados com Docker Engine 29.1.3 e Compose 2.40.3, em Ubuntu/WSL2. O host informa processador AMD Ryzen 7 5825U, 16 CPUs lógicas disponíveis e aproximadamente 15,3 GiB de memória na máquina virtual. Os manifestos preservam versões efetivas, identificadores de imagens e configuração de recursos. A imagem por tag não é tratada como identificação imutável; o digest registrado permite distinguir o artefato usado.

\subsection{Arquitetura de software}
Os serviços permanentes são \texttt{postgres}, \texttt{kafka}, \texttt{sync-api}, \texttt{async-api} e \texttt{consumer}. Serviços transitórios preparam permissões do volume Kafka e criam explicitamente os tópicos. A API síncrona executa validação, hash e persistência antes de responder. A assíncrona aguarda o ACK individual do broker; o consumidor realiza o processamento posterior com o mesmo repositório.

O projeto Compose dedicado utiliza as portas locais 18000 e 18001 para as APIs, 15432 para PostgreSQL e 19092 para Kafka, vinculadas à interface de loopback. Os endereços internos são próprios da rede Compose. Volumes nomeados preservam banco e log Kafka; as políticas de reinício permitem retomar os processos. Não se utilizam dados ou serviços de produção.

Cada aplicação possui limite de uma CPU e 256 MiB; PostgreSQL, uma CPU e 512 MiB; Kafka, uma CPU e 1 GiB. Esses limites delimitam o ambiente de desenvolvimento, mas ainda requerem avaliação no piloto. Os tópicos \texttt{audit-events} e \texttt{audit-events-dlq} possuem uma partição e fator de replicação um. A topologia não demonstra alta disponibilidade nem tolerância à perda física do único broker.

\subsection{Estrutura dos módulos}
A Tabela~\ref{tab:modulos-implementados} relaciona os arquivos existentes às responsabilidades. O consumidor utiliza funções de módulo; os diagramas não pressupõem classes adicionais ausentes do código.

{\small
\begin{longtable}{p{0.38\textwidth}p{0.53\textwidth}}
\caption{Módulos implementados e responsabilidades}\label{tab:modulos-implementados}\\
\hline Caminho em implementacao/ & Responsabilidade\\ \hline
\endfirsthead
\hline Caminho em implementacao/ & Responsabilidade\\ \hline
\endhead
\texttt{app/sync\_api/main.py} & POST síncrono, consulta por UUID, prontidão e vitalidade.\\ \hline
\texttt{app/async\_api/main.py} & Validação, publicação confirmada e resposta 202; tratamento de resultado incerto.\\ \hline
\texttt{app/consumer/main.py} & Consumo, validação, retry, DLQ, confirmação de posição e retomada da sessão.\\ \hline
\texttt{app/shared/schemas.py} & Contratos de evento, aceite, persistência e consulta.\\ \hline
\texttt{app/shared/processing.py} & Serialização determinística e SHA-256.\\ \hline
\texttt{app/shared/database.py} & Esquema, transação, unicidade concorrente, conflito e consulta.\\ \hline
\texttt{app/shared/kafka.py} & Produtor com callback por mensagem, prazo de confirmação e recibo de entrega.\\ \hline
\texttt{app/shared/observability.py} & Logs JSONL com campos permitidos, UUID de execução e marcos temporais.\\ \hline
\texttt{tests/} & Testes unitários, integração e verificações de falhas controladas.\\ \hline
\texttt{scripts/load.js} & Gerador k6 com chegada aberta, identificação de fase e correlação de eventos.\\ \hline
\texttt{scripts/run\_experiment.py} & Orquestração, manifestos, coleta e bloqueio do protocolo incompleto.\\ \hline
\texttt{scripts/normalize.py} & Conciliação de eventos, métricas descritivas e identificação de inconsistências.\\ \hline
\texttt{scripts/plot\_validation.py} & Gráfico de validação, tabela e hashes das fontes usadas.\\ \hline
\texttt{docker-compose.yml} & Topologia isolada, recursos, volumes, dependências e testes de prontidão.\\ \hline
\multicolumn{2}{l}{\footnotesize Fonte: Elaborado pelo autor, a partir do código-fonte (2026).}
\end{longtable}
}

\subsection{Funcionalidades implementadas e limites}
A persistência usa \texttt{INSERT ON CONFLICT DO NOTHING RETURNING} e chave primária no UUID. Em caso de duplicata, uma nova consulta sob isolamento \texttt{READ COMMITTED} verifica o hash. Conteúdo idêntico é reconhecido sem nova linha; conteúdo divergente produz conflito e preserva o registro original. A resposta síncrona sucede o término da transação. O campo \texttt{persisted\_at} recebe o horário de inserção; o log \texttt{commit\_completed} é emitido depois do commit.

O produtor habilita idempotência da biblioteca e \texttt{acks=all}, aguarda o callback individual e fornece tópico, partição e offset. Buffer indisponível não gera aceite; timeout e erros de entrega podem gerar resposta 503 com aceitação desconhecida. Nessa situação, não se presume que a mensagem tenha sido perdida, e um reenvio deve preservar o UUID. A idempotência do produtor não constitui transação distribuída com PostgreSQL.

O consumidor exige UUID e instante de ocorrência explícitos na mensagem Kafka, evitando gerar nova identidade durante reentregas. A política local prevê três tentativas totais de persistência, com esperas de 0,5 e 1 segundo na configuração usada. Falhas terminais, mensagens inválidas e conflitos são retidos na DLQ com motivo e posição de origem. O offset de origem somente avança após persistência ou ACK da DLQ. Se essa confirmação falhar, a sessão é retomada com a posição não resolvida.

A DLQ preserva os bytes originais em Base64, diferenciando valor nulo de conteúdo vazio. Sua mensagem inclui uma identidade formada por tópico, partição e offset de origem. Uma falha entre o ACK da DLQ e o commit da origem pode produzir nova cópia na DLQ; não se promete entrega exatamente uma vez. A retomada do serviço também não implica reprocessamento automático das mensagens já encaminhadas à DLQ.

As APIs distinguem vitalidade de prontidão: um processo pode estar ativo e retornar 503 no teste de dependências. A consulta por UUID retorna uma projeção do registro, sem expor o payload. Os logs correlacionáveis utilizam lista permitida de campos e não exportam credenciais. O painel dedicado, as quatro telas desenhadas e as tabelas complementares propostas não estão implementados; a operação atual ocorre por API, scripts e arquivos.

\subsection{Testes e validação funcional}
\label{sec:validacao-funcional}
A Tabela~\ref{tab:testes-executados} resume execuções reais. Os testes unitários com doubles verificam decisões locais; a integração e os ensaios de falha utilizam PostgreSQL e Kafka em contêineres. As 22 verificações pertencem a uma sequência funcional controlada, e não a 22 repetições estatísticas.

{\small
\begin{longtable}{p{0.24\textwidth}p{0.12\textwidth}p{0.55\textwidth}}
\caption{Verificações executadas no protótipo}\label{tab:testes-executados}\\
\hline Grupo & Aprovados & Cobertura observada\\ \hline
\endfirsthead
\hline Grupo & Aprovados & Cobertura observada\\ \hline
\endhead
Unitários da aplicação & 24 & Contrato, hash, ACK, erros, retry, DLQ e decisões de confirmação.\\ \hline
Integração & 5 & Oito inserções concorrentes do mesmo UUID, conflito, publicação/consumo, mensagem inválida e conflito assíncrono.\\ \hline
Falhas controladas & 22 & Parada e retomada do consumidor, Kafka e banco; conciliação de UUIDs, DLQ e offsets.\\ \hline
Instrumentos offline & 36 & Fórmulas, fases, ausência de dados, deduplicação, arquivos, hashes e bloqueios de elegibilidade.\\ \hline
\multicolumn{3}{l}{\footnotesize Fonte: Relatórios de execução e arquivos de evidência do autor (2026).}
\end{longtable}
}

No teste de concorrência, oito submissões produziram uma inserção e sete reconhecimentos de duplicata, com uma linha no banco. Com o consumidor parado, três eventos receberam 202 e permaneceram ausentes do PostgreSQL; depois da retomada, os três foram persistidos e o lag observado passou de três para zero.

Com Kafka indisponível, foi observada resposta 503 com aceitação desconhecida, e não 202. Com PostgreSQL indisponível, a API síncrona sinalizou falta de prontidão, enquanto um evento aceito pela assíncrona esgotou as três tentativas e foi retido na DLQ. Após a volta do banco, um novo evento foi persistido sem reinício manual do consumidor. O evento já encaminhado à DLQ não foi indevidamente classificado como persistido.

\subsection{Evidências e capturas de execução}
\label{sec:evidencias}
As Figuras~\ref{fig:evidencia-api} a~\ref{fig:evidencia-recuperacao} são capturas de um relatório local produzido a partir dos registros brutos. Não são telas de um painel de produção nem saídas de terminal reconstruídas. A apresentação seleciona campos para legibilidade; os arquivos originais, UUID da execução e hashes permanecem disponíveis no diretório de revisão. As evidências de falha pertencem à execução \texttt{b4e7c2a8-a36f-4c90-b481-7250d70a20d8}.

\newcommand{\evidenciafig}[3]{%
\begin{figure}[htbp]
\centering
\includegraphics[trim=0 55bp 0 0,clip,width=0.98\textwidth,height=0.42\textheight,keepaspectratio]{Imagens/#1.png}
\caption{#2}\label{#3}
\fonte{Captura do relatório de execução, elaborado pelo autor com dados reais de teste (2026).}
\end{figure}
}

A Figura~\ref{fig:evidencia-api} apresenta requisição sintética e resposta assíncrona; a Figura~\ref{fig:evidencia-publicacao} identifica o ACK que antecedeu essa resposta. A Figura~\ref{fig:evidencia-consumo} correlaciona o mesmo evento ao marco pós-commit e à confirmação da posição consumida.

\evidenciafig{evidencia-api}{Requisição real e aceite assíncrono durante a parada do consumidor}{fig:evidencia-api}
\evidenciafig{evidencia-publicacao}{Confirmação individual de publicação no Kafka}{fig:evidencia-publicacao}
\evidenciafig{evidencia-consumo}{Persistência e confirmação posterior do offset}{fig:evidencia-consumo}

A Figura~\ref{fig:evidencia-banco} reúne a linha exportada por SQL e as respostas a reenvios idêntico e conflitante. As Figuras~\ref{fig:evidencia-falhas} e~\ref{fig:evidencia-recuperacao} documentam tentativas, retenção e retomada. Retenção na DLQ não equivale a conclusão do processamento de auditoria.

\evidenciafig{evidencia-banco}{Consulta PostgreSQL, duplicata e conflito verificados}{fig:evidencia-banco}
\evidenciafig{evidencia-falhas}{Tentativas limitadas e retenção de falha terminal na DLQ}{fig:evidencia-falhas}
\evidenciafig{evidencia-recuperacao}{Síntese das interrupções e recuperações observadas}{fig:evidencia-recuperacao}
\FloatBarrier

\subsection{Verificação dos instrumentos de coleta}
A coleta curta utilizou k6 com cinco eventos por segundo, durante dez segundos, sem aquecimento, em cada variante. Esses parâmetros foram escolhidos exclusivamente para exercitar os instrumentos. Não representam a taxa de referência do experimento nem uma calibração de capacidade. Cada execução enviou, aceitou e conciliou cinquenta eventos, sem erros HTTP ou iterações descartadas observados.

A Figura~\ref{fig:evidencia-metricas} apresenta a saída consolidada real. As medianas de conclusão foram 6,996 ms e 15,433 ms, respectivamente, nas execuções síncrona e assíncrona. Esses números descrevem somente duas coletas curtas, uma por variante, e não permitem afirmar superioridade arquitetural. Ambas utilizaram o mesmo marco pós-commit. As execuções foram identificadas por \texttt{e029f683-468e-4235-a474-9c2e6ca8213d} e \texttt{ac66bb9a-aa21-4210-ba19-8f5a90dea42b}.

\evidenciafig{evidencia-metricas}{Métricas reais para validação do instrumento, sem inferência comparativa}{fig:evidencia-metricas}

Cada diretório contém eventos preparados, marcos do gerador, saída k6, logs JSONL, exportação SQL, recursos, lag, configuração, manifesto, snapshot dos fontes e checksums. O normalizador preserva ausências como dados ausentes, não como zero, e separa aceite, commit, erro, DLQ e pendência. Estatísticas descritivas não são rotuladas como conclusão científica.

A consulta de lag inicia uma JVM pela ferramenta Kafka dentro do contêiner do broker. Seu consumo interfere na observação de CPU e memória; o piloto deverá quantificar esse custo ou selecionar um coletor menos intrusivo antes do congelamento. Os resumos de recursos da validação abrangem toda a observação, inclusive drenagem, e não sustentam conclusões de eficiência.

\subsection{Reprodutibilidade e etapas restantes}
Os comandos completos estão nos arquivos \texttt{README.md} da implementação e dos scripts. A aplicação pode ser validada por testes, e a coleta curta pode ser repetida pelo comando abaixo, em Linux/WSL com a topologia isolada saudável:
\begin{Verbatim}[fontsize=\small,breaklines=true,frame=single]
python3 scripts/run_experiment.py validation \
  --variant both --duration 10 --rate 5 --drain 10
\end{Verbatim}

O protocolo definitivo permanece deliberadamente não congelado. O executor recusa configuração incompleta; exige registro do piloto, definição de C4, parâmetros de carga e falha, 60 segundos de aquecimento, 300 segundos de medição e pelo menos dez repetições. O congelamento depende da calibração e da decisão metodológica discutida no Capítulo~\ref{cap:metodo}. A coleta científica, sua análise e as conclusões comparativas ainda não foram realizadas.

\textbf{Síntese.} As duas variantes e seus mecanismos de confirmação, idempotência, recuperação e coleta foram verificados funcionalmente. As evidências são rastreáveis aos dados brutos. A pergunta de pesquisa ainda exige piloto e repetições controladas.
